\cotefile{84}
\scanpage{130}

\[
\begin{aligned}
\unsure{0}\to L_i - \deleted{\unsure{F_i}}E_i - M_i - 0
 &,\ \unsure{\text{et}}\ \pi_i:E_i-L_i\\
 &\qquad\unsure{\chi|\lambda}\text{-scindage}\\
p_i:E_i-L_i &\qquad\unsure{\text{à gauche}}\qquad p=(p_1,p_2)\\[-2pt]
\text{scindage}&
\end{aligned}
\]

\[
\begin{aligned}
\Theta_{p_i}&:E_i\simeq M_i\oplus L_i\\
b_i&:\pi_i-\chi p_i:M_i\longrightarrow L_i\\
b&=b_1\otimes b_2:\deleted{\unsure{L}}M_1\otimes M_2\longrightarrow L_1\otimes L_2
\end{aligned}
\]

\[
\begin{aligned}
p'_i&=p_i-u_i,\quad u_i:M_i\to L_i\\
p'&=(p'_1,p'_2)\\
\text{d'où }&b'_1,b'_2,\ \text{avec}\\
\deleted{\unsure{u}}b'_1&=\deleted{\unsure{u}}b_1+\chi u_1,\quad b'_2=b_2+\chi u_2
\end{aligned}
\]

d'où $\chi$-trivialisation \unsure{complète}, \unsure{structure|extension}$_{(p_1,p_2)}$

\diagram{\begin{tikzcd}[ampersand replacement=\&,column sep=large]
\mathcal E_{(p_1,p_2)}\deleted{\unsure{M_1\otimes M_2\oplus{}}}=(M_1\otimes M_2)\oplus(L_1\otimes L_2)
 \arrow[r,"{\deleted{\unsure{\Theta}}}","{\Theta_{p,p'}}"'] \&
 \mathcal E_{p'_1,p'_{2}}
\end{tikzcd}}

\[
\deleted{\unsure{\Theta_{p,p'}}}\quad\deleted{={}\,T_{-}}u
\]

\[
\begin{aligned}
\pi_{(p_1,p_2)}-\chi p&=b_1\otimes b_2\\
\pi_{(p_1,p_2)}:(\unsure{y|\eta},\unsure{z|\xi})&\longrightarrow
 \deleted{\unsure{z}}\chi z+b_1\otimes b_2(y))
\end{aligned}
\]

\[
(\unsure{y|\eta},\unsure{z|\xi})\mapsto(y,z\unsure{-|+}uy)
\]

\[
u=b_1\otimes u_2+u_1\otimes b_2+\chi u_1\otimes u_2
\]

\margin{\uline{NB}\quad $\Theta_{p,p'}$ transforme $\pi_{(p_1,p_2)}$}

\[
\begin{aligned}
\pi_{p'_1,p'_2}:(y,z)&\longmapsto\deleted{\unsure{z}}\chi z+(b'_1\otimes b'_2)(y)\\
 &=\pi_{(p_1,p_{\unsure{1|2}})}(y,z)+u(y)
\end{aligned}
\]

% archive sheet number: 129
\scanpage{131}

\mspage{468}

(\textsuperscript{107}) (4 octobre) J'ai eu occasion déjà de faire mention d'un aspect important de ces premières cinq années de ma vie, comme d'un"privilège" de grand prix (*): une identification profonde et sans problèmes avec mon père, laquelle n'a jamais été touchée par la peur ou par l'envie. Je me suis rendu compte de cette circonstance, et de l'existence même, comme de la silencieuse force, de cette identification à mon père, il y a quatre ans seulement (au cours de la méditation sur mon enfance et sur ma vie qui a suivi celle d'août 79 à mars 80 sur mes parents) Cette identification était comme le coeur paisible et puissant d'une identification à la famille que nous formions, mes parents, ma soeur (qui était mon \unsure{aînée|aîné} de quatre ans) et moi. Je vouais une admiration et un amour sans limites à mon père comme à ma mère. Leur personne était pour moi la mesure de toutes choses.

Cela ne signifie nullement que mon attitude à leur égard était celle \unsure{d'une} approbation automatique, d'une admiration béate. Je ne savais sans doute pas qu'ils étaient la mesure de toute chose pour moi, mais je savais fort bien qu'il étaient faillibles comme moi, et il n'y avait en moi aucune crainte qui m'aurait empêché de constater un désaccord et de le manifester clairement\unsure{..|.} Dans les conflits qui m'entouraient\unsure{;|,} je ne craignais pas de prendre partie à ma façon. Ce ne touchait en rien à une certaine foi, à une assurance qui formaient l'assise profonde, inébranlable de mon être - plutôt, cela découlait spontanément de \unsure{cette} foi, de cette assurance même.

Il arrivait que mon père, dans des accès de colère impuissante alors que ma soeur (sans en avoir l'air) prenait plaisir à le provoquer, la frappe avec brutalité - et à chaque fois j'en étais outragé, dans un élan de solidarité sans réserves avec ma soeur. C'étaient là je crois les seuls gros nuages qui \unsure{passaien|passaient} dans ma relation à\unsure{:}mon père (il n'y en avait pas avec ma mère). Ce n'est pas que j'approuvais les tours parfois pendables de ma soeur, pas plus je crois qu'ils ne me troublaient vraiment - ce n'était pas \mbox{e l l e} qui était pour moi la mesure des choses. Ses tours (dont la raison sûrement m'échappait tout autant qu'à mon père, qui "marchait" à tous les coups, ou à ma mère qui n'avait garde d'intervenir ni avant ni après) - ces tours en un sens ne tiraient pas vraiment à conséquence pour moi. C'était ma soeur, elle était comme elle était, voilà \unsure{tou|tout} Mais que \mbox{m o n}\quad\mbox{p è r e} se laisse aller à une telle brutalité aveugle...

Les trois êtres les plus proches, qui ensemble ont constitué comme la matrice de mes premières années, étaient déchirés par le conflit, opposant chacun

\noindent\rule{5em}{0.4pt}\par
(*) Voir la note "Le massacre", n° 87.

% archive sheet number: 130
\scanpage{132}

\[
\begin{aligned}
\deleted{\Theta_{p+u}(x)=}\qquad
\Theta_{p}(x)&=\deleted{\unsure{p(x)}}\unsure{\psi|\varphi}(x)\oplus p(x)\\
\Theta_{p+u}(x)&:\psi(x)\oplus(p(x)+u\psi(x))
\end{aligned}
\]

\margin{$\boxed{b_i=(\pi_{\unsure{1|i}}-\unsure{\chi|\lambda}p_i)}$}

\diagram{\begin{tikzcd}[ampersand replacement=\&,column sep=large,row sep=huge]
 \& E_1\mathbin{\unsure{\vee|\cup}} E_2
 \arrow[dl,"{\Theta_{p_1\ast p_2}}"',"\sim",sloped]
 \arrow[d,"{\Theta_{p'_1\ast p'_2}}"',"\wr"]
 \arrow[dr,"{\Theta_{p''_1\ast p''_{\unsure{3|2}}}}"',"\sim",sloped]\&\\
 M\oplus L\arrow[r,"{\overset{(u_1,u_2,u)}{\Theta_{(p_1,p_2)(p'_1,p'_2)}}}","\sim" description]
 \arrow[rr,bend right=24,"{u''_1,u''_2,u''}","{\Theta_{(p_1,p_2)(p''_1,p''_2)}}"']\&
 M\oplus L\arrow[r,"{\overset{(u'_1,u'_2,u')}{\Theta_{(p'_1,p'_2)(p''_1,p''_{\unsure{3|2}})}}}","\sim" description]\&M\oplus L\\
 \operatorname{Hom}(M_i,L_i)\arrow[d,phantom,"\rotatebox{90}{$\in$}"]\&
 \operatorname{Hom}(M_{\unsure{2|i}},L_i)\arrow[d,phantom,"\rotatebox{90}{$\in$}"]\&
 \operatorname{Hom}(M_i,L_i)\\
 b_i,\& b'_i=b_i+\chi u_i\& b''_i=b'_i-\chi\deleted{\unsure{b}}u'_i
\end{tikzcd}}

\[
\begin{aligned}
(y,z)&\mapsto (y,z-uy)\qquad (y,z)\mapsto(y,z-u'y)\\
(y,z)&-(y,z+
\end{aligned}
\]

\[
\left\{\begin{aligned}
p'_1\ast p'_2&=p_1\ast p_2-u\\
p''_1\ast p''_2&=p'_1\ast p'_2-u'=p_1^{\unsure{\ast}}\ast p_2-u''\\
&\Longrightarrow u''=u+u'
\end{aligned}\right.
\]

\[
\begin{aligned}
 b_1\otimes b_2&\in\operatorname{Hom}(M,L)& b'_1\otimes b'_2&\\
 (\pi-\chi\,p_1\ast p_2)^{\Theta_{p_1\ast p_2}}& &
 &=(b_1-\chi u_1)\otimes(b_2-\chi u_2)\\
 &&&=b_1\otimes b_2\unsure{-|+}\chi u
\end{aligned}
\]

\[
b''_1\otimes b''_2=b'_1\otimes b'_2-\chi u'
=b_1\otimes b_2-\chi\underbrace{(u+u')}_{u''}
\]

\diagram{\begin{tikzcd}[ampersand replacement=\&,row sep=tiny]
(p_1,p_2)\arrow[d]\\(p'_1,p'_2)\arrow[d]\\
\underline{(p''_1,p''_2)}\\
(p_1,p_2)\arrow[d]\\(p''_1,p''_2)
\end{tikzcd}}

\[
\left\{\begin{aligned}
u&=b_1\otimes u_2+u_1\otimes b_2+\chi u_1\otimes u_2\\
u'&=b'_1\otimes u'_2+u'_{\unsure{1}}\otimes b'_2+\chi u'_1\otimes u'_2\\
u''&=b_1\otimes u''_2+u''_1\otimes b_2+\chi u''_1\otimes u''_2
\end{aligned}\right.
\]
(où $u''_1=u_1+u'_1,\ u''_2=u_2^{\deleted{\unsure{\prime}}}+u'_2$)

\[
\begin{aligned}
&\underbrace{b_1,b_2\,;\,\underbrace{u_1,u_2}_{\downarrow}} \qquad
 \underbrace{u'_1,u'_2}_{\unsure{\vert|\mathrm{I}}}\\
&\deleted{\unsure{u=}}\quad
\left\{\begin{aligned}
 b'_1&=b_1+\chi u_1,\ b'_2=b_2+\chi u_2\\
 b''_1&=b'_1+\chi u'_1,\ b''_2=b_2+\chi u'_2
\end{aligned}\right.
\quad\left|\begin{aligned}
 b&=b_1\otimes b_2\\ b'&=b'_1\otimes b'_2\\ b''&=b''_1\otimes b''_2
\end{aligned}\right.\\
&\deleted{(b'_1-\chi u_1)\otimes u'_1+u'_2\otimes(b_2-\chi u_2)+\chi u'_1\otimes\unsure{u'_2}}\\
&\deleted{(b'_1,b'_2)}
\end{aligned}
\]

1°) Je dis que $u''=u+u'$, \deleted{$b''_1=b'_1-\chi u''_1,\ b''_2=\unsure{b_2}=$}

\[
\begin{aligned}
u''&=b_1\otimes\unsure{(}\underset{\cdot}{u}_2+u'_2)+(\underset{\cdot}{u}_1^{\deleted{\unsure{\prime}}}+u'_1)\otimes \underset{\cdot}{b}_2+\chi\underline{(u_1+u_{\unsure{2|1}}')\otimes{}}(u_2+u'_2).\\
&=u+b_1\underset{\circ}{\otimes} u'_2+
\underbrace{u'_1\otimes b_2+\chi u'_1\otimes u_2}_{u'_1\otimes(b_2+\chi u_2)}
+\chi\underset{\cdot}{u}_1\otimes u'_2+\chi u'_1\otimes u'_2\\
&\hspace{8em}+(b_1+\chi u_1)\otimes u'_2+\chi u'_1\otimes u'_2
\end{aligned}
\]
\margin{$=\deleted{\unsure{u}}$}

% archive sheet number: 131
\scanpage{133}

\mspage{469}

d'eux et à lui-même, et aux deux autres: conflit insidieux,au visage impassible entre ma mère et ma soeur, et confli\unsure{t} aux violents éclats entre mon père et ma mère d'un côté, ma soeur de l'autre, qui chacune pour son propre compte (et sans que personne du vivant de mes parents ait jamais fait mine de s'en apercevoir... le faisait marcher à sa foçon. La chose mystérieuse, extraordinaire, c'est qu'entouré ainsi par le conflit en ces années les plus sensibles, les plus cruciales de la vie, celui-ci soit resté \mbox{e x t é r i e u r} à moi, qu'il n'ait pas vraiment "mordu" sur mon être en ces années-là et ne s'y soit installé à demeure.

La division dans mon être, qui a marqué ma vie tout autant que celle de \unsure{to|tout} autre, ne s'est pas installée en moi en ces années, mais en les deux ou trois \unsure{an|ans} qui ont suivi, de ma sixième à ma huitième année environ. A un certain moment. (que j'ai cru pouvoir situer à quelques mois près, et qui se placerait dans ma huitième année) il y a eu un certain \mbox{b a s c u l e m e n t}, au bout de plus de deux ans de séparation avec mes parents (qui ne se souciaient guère de me donner signe de vie) et de ma soeur. C'était avant tout une \mbox{r u p t u r e}\quad\mbox{a v e c}\quad\mbox{m o n}\quad\mbox{e n f a n c e}, "enterrée" à partir de ce moment par d'efficaces mécanismes d'oubli (qui sont restés en place, à peu de choses près, jusqu'à aujourd'hui même). A un certain niveau profond (pas le plus profond pourtant...) mes parents ont alors été déclarés par moi comme des "étrangers", tout comme mon enfance était désormais déclarée "étrangère". J'ai \mbox{a b\unsure{.}d i q u é}, en un sens: pour être accepté dans le monde qui m'entourait désormais, j'ai décidé d'être comme "eux", comme les adultes qui y font la loi - d'acquérir et de développer les armes qui y forcent le respect, de me battre à armes égales dans un monde où seule, une certaine espèce de "force" est acceptée et prisée...

C'était d'ailleurs cette force-là aussi qui avait la préférence de mes parents, lesquels avaient entouré mes premières années. Et là je reviens à cette "chose mystérieuse" (dont je viens de m'éloigner, en suivant le fil d'une autre association suscitée par cette chose), \mbox{l' a b s e n c e}\quad\mbox{d e}\quad\mbox{d i v i s i o n}\quad\mbox{e n}\quad\mbox{m o i}, en ces premières années de ma vie.

Peut-être le mystère n'est plus pour moi en cette absence, mais plutôt en ceci: que mes parents, mon père comme ma mère, m'aient chacun alors \mbox{\unsure{a c c e p t|a c c e p t é}}\quad\mbox{d a n s}\quad\mbox{m a}\quad\mbox{t o t a l i t é}, et totalement: dans ce qui en moi est "viril" est "homme", et dans ce qui est "femme". Ou pour le dire autrement: que mes parents, déchirés l'un et l'autre par le conflit, reniant chacun une partie essentielle de leur être - incapable chacun d'une ouverture aimante à lui-même et à l'autre, comme d'une ouverture aimante à ma soeur... que néanmoins ils aient trouvé une telle ouverture, une acceptation sans réserve, vis à vis de moi

% archive sheet number: 132
\scanpage{134}

\diagram{\begin{tikzcd}[ampersand replacement=\&,column sep=large,row sep=large]
\& E_1\mathbin{\unsure{\vee|\ast}}E_2\arrow[dl,"{\Theta_{p_1\ast p_2}}"']\arrow[dr,"{\Theta_{p'_1\ast p'_2}=\Theta_{p_1\ast p_2}\unsure{-|+}u}"]\&\\
(M_1\otimes M_2)\oplus(L_1\otimes L_2)
\arrow[rr,"{\unsure{\Lambda|\lambda}_u:(\deleted{\unsure{z}}y,\deleted{\unsure{y}}z)\longmapsto(y,z+u(y))}"'] \& \&
(M_1\otimes M_2)\oplus(L_1\otimes L_2)\\
b_1,b_2\&\& b'_1,b'_2
\end{tikzcd}}

\[
\begin{aligned}
u&=b_1(p_1)\otimes u_2+u_1\otimes b_2(p_2)\\
 &\qquad+\unsure{\chi|\lambda}u_1\otimes u_2
\end{aligned}
\]

\[
b'_1=b_1-\chi u_1,\qquad b'_2=b_2-\chi u_{\unsure{2}}
\]

\[
\begin{aligned}
p'_1&=p_1\unsure{-|+}u_1\unsure{\psi|\varphi}_1,\ p'_2=p_2\unsure{-|+}u_2\psi_2,\\
u_1&:M_1\to L_1\\ u_2&:M_2-L_{\unsure{2}}
\end{aligned}
\]

\[
p'_1\ast p'_2=p_1\ast p_2\unsure{-|+}
\left[\underbrace{(\pi_1-\chi p_1)}_{b_1(p_1)}\otimes u_2+
u_1\otimes\underbrace{(\pi_2-\chi p_2)}_{b_2(p_2)}\right]
\deleted{\unsure{+}}-\chi u_1\otimes u_2
\]
\deleted{\unsure{avec}}

\margin{$\uparrow$ signe étrange qui \unsure{doit|est} \unsure{également|prédiction} \unsure{être} \deleted{\unsure{étrange}}}

\[
u_1\ast u_2=-\chi u_1\otimes u_2
\]

Ces signes sont \unsure{dûs|dus} au signe $\underset{\unsure{\varepsilon}}{-}\mathrm{id}$ dans le \unsure{diagr.|diag.} \unsure{non} commutatif. (pour une extension \unsure{de} $\underline O$ par $\mathcal H$

\diagram{\begin{tikzcd}[ampersand replacement=\&,column sep=small]
0\arrow[r,dash]\&\mathcal H\arrow[r,"{\alpha}",dash]\arrow[d,"{-\mathrm{id}}"]\&
\mathcal E\arrow[r,"{\psi}",dash]\arrow[d,"{\unsure{\mathcal D|\Phi}}","{\unsure{\Phi}}"']\&
\underline O\arrow[r]\arrow[d,equal]\&0\\
0\arrow[r,dash]\&\mathcal H\arrow[r]\&\operatorname{Hom}^{\unsure{\prime|!}}(E,\mathcal H)\arrow[r]\&O\arrow[r,dash]\&0\\
\&\operatorname{Hom}(\underline O,\mathcal H)\arrow[u,phantom,"\unsure{\Uparrow}"]\&\&\&
\end{tikzcd}}

\[
\unsure{\Phi|\mathcal D}(x)(y)=-x\psi(y)+y\psi(x)\qquad\unsure{\simeq}
\]

\[
\begin{aligned}
&\deleted{=\{f:\mathcal E\to\mathcal H\quad f\mapsto\lambda}\\
&\deleted{\text{\unsure{qui a la vert.}}\ \exists\lambda\,;\,}\\
&=\{f:\mathcal E-\mathcal H\mid f\circ\alpha=\alpha(\lambda\mathrm{id}_{\mathcal H}\}
\end{aligned}
\]

\diagram{\begin{tikzcd}[ampersand replacement=\&,row sep=tiny]
\&\arrow[r]\&\deleted{\unsure{z}}\\
\unsure{z|\xi}_1,\ y_1\arrow[rr]\&\&\chi y_1+b_1(z_1)\\
{M_1\otimes M_2,\quad L_1\otimes L_2}\arrow[u,phantom,"{\cap\quad\cap}"]\&\&
\end{tikzcd}}

Si on pose $p'_1=p_1-u'_1,\ p'_2=p_2-u'_2,\ p'_1\ast p'_2=p_1\ast p_2-u'$

\[
\begin{aligned}
\text{d'où }\quad&\left\{\begin{aligned}u'_1&=-u_1\\u'_2&=-u_2\\u'&=-\unsure{u|u_3}\end{aligned}\right.
\qquad\text{on aura}\quad u'=-b_1\otimes u_2-u_1\otimes b_2+\chi u_1\otimes u_2\\
&\boxed{u'=b_1\otimes u'_2+u'_1\otimes b_2+\chi u'_1\otimes u'_2}
\end{aligned}
\]

\margin{c'est la \unsure{formule primitive|formule transitive}....}

% archive sheet number: 133
\scanpage{135}

\mspage{470}

leur fils.

Pour le dire autrement encore: à aucun moment en ces premières\unsure{,} cinq années de ma vie, je n'ai connu le sentiment de \mbox{h o n t e}\quad\mbox{d' ê t r e}\quad\mbox{c e}\quad\mbox{q u e}\quad\mbox{j e}\quad\mbox{s u i s}, que ce soit dans mon corps et ses fonctions, ou dans mes \unsure{pulsion|pulsions} mes penchants,\unsure{.} mes actions. A aucun moment je n'ai eu à renier quelque chose en moi, pour être accepté par mon entourage et pouvoir vivre en paix avec lui.

Il arrivait bien sûr que je fasse des choses qui ne "passaient" pas: comme tous les enfants il m'arrivait sûrement d'être pénible, voire insupportable quand je m'y mettais - et il était clair\unsure{.} parfois qu'il me fallait rectifier le tir. Je ne faisais pas la loi, ni n'étais tenté de vouloir la faire, n'ayant pas à compenser quelque mutilation secrète. Et dans l'amour de mes parents pour moi, il n'aurait pu y avoir de place pour l'adulation, la complaisance aux caprices - pour une approbation inconditio\deleted{\unsure{ne}}nne\unsure{l}le. Mais s'il arrivait forcément que je me fasse "envoyer sur les \deleted{\unsure{c}}roses" par mon père ou par ma mère (tout comme l'inverse pouvait parfois se produire), jamais dans ces années-là l'un ni l'autre ne m'ont fait honte, d'un acte ou d'un comportement qui n'aurait pas eu l'heur de leur plaire.

Sur le fond d'une identification profonde au père, sans ambiguité aucune, ma personne comme enfant m'apparaît aujourd'hui comme empreinte à la fois de virilité et de féminité, fortes l'une comme l'autre.

Il me semble qu'en chaque être et en chaque chose, dans ces indissolubles et fluctuantes épousailles des qualités yin et yang en lui qui font de lui ce qu'il est, et dont le délicat équilibre est la beauté profonde, l'harmonie qui vit en cet être ou en cette chose - que dans cette union intime du yin et du yang il y a souvent (peut-être toujours) une note de fond, une "dominante", qui est soit yin, soit yang. Cette note de fond n'est pas toujours aisée à déceler en une personne, à cause des mécanismes de répression plus ou moins efficaces et complets, qui faussent le jeu en substituant, à une harmonie originelle une image d'emprunt. Ainsi mon "image de marque" pendant quarante ans a été une image presque exclusivement virile - sans d'ailleurs jamais se voir mettre en cause ni même être décelée comme telle, par moi-même ni (il me semble) par autrui, jusqu'à ma quarante-huitième année. J'ai tendance à croire pourtant que la note de fond présente à la naissance reste présente pendant la vie entière, dans des couches profondes tout au moins qui jamais peut-être ne trouveront l'occasion de se manifester au grand jour. Dans mon propre cas, chose étrange, \unsure{j|je} ne saurais aujourd'hui encore dire quelle est cette note dominante, celle donc

% archive sheet number: 134
\scanpage{136}

\diagram{\begin{tikzcd}[ampersand replacement=\&,column sep=small]
0\arrow[r]\&L\arrow[r,dash]\arrow[d,"{-\mathrm{id}_L}"]\&
 E\arrow[r,"{\unsure{\psi|\varphi}}"]\arrow[d,"{\wr}"]\&
 \underline O\arrow[r]\arrow[d,equal]\&0\\
 {}\arrow[r]\&L\arrow[r]\&\operatorname{Hom}^{!}(E,L)\arrow[r]\&
 \underline O\arrow[r,dash]\&0
\end{tikzcd}}

\[
\begin{aligned}
\unsure{x}&\longmapsto[y\longmapsto(-x\psi(y)+y\psi(x))]\\
&E\xrightarrow[\psi(x)]{}L
\end{aligned}
\]

\[
{}_{\lambda_1}E_1\ast{}_{\lambda_2}E_2\simeq
\]

\[
\boxed{\begin{aligned}
(\deleted{\unsure{L}}f_1\ast f_2)(x_1\ast x_2)
 &=\unsure{\psi_1}(f_1)\pi_1(x_1)\otimes f_2(x_2)
 + f_1(x_2)\otimes\psi_2(f_2)\pi_2(x_2)
 -\unsure{\chi|\lambda} f_1(x_1)\otimes f_2(x_2)\\
(f_1\ast f_2)(u_1\otimes u_2)&=\lambda_1\lambda_2\,u_1\otimes u_2
\end{aligned}}
\]

\[
(f_1\ast u_2)(x_1\ast x_2)=
\psi_1(f_1)\pi_1(x_1)\otimes\underbrace{\deleted{\unsure{f}}u_2(x_2)}_{u_2(\psi_2(x_2))}
-\chi f_1(x_1)\otimes\deleted{\unsure{f}}u_2(x_2)
\]

\deleted{$(u_1\ast f_2)(x_1\ast x_2)$}

\[
\begin{aligned}
\text{i.e.}\quad (f_1\ast u_2)(x_1\ast x_2)
 &=\underbrace{\big(\psi_1(f_1)\pi_1-\overset{\substack{\psi_1(\pi_1)\\ \shortparallel}}{\chi} f_1\big)}_{\text{nul sur } L_1}\otimes\deleted{\unsure{f}}u_2(x_1\otimes x_2)\\
 &\qquad\big([\pi'_1,f'_1]'\otimes u_2\big)(x_1\otimes x_2)
\end{aligned}
\]

\[
\text{\unsure{de|du|donc|d'où}}\quad =\quad\left\{\begin{aligned}
u_1\ast f_2\,(x_1\ast x_2)&=(u_1\otimes[\pi'_2,f'_2]')(x_1\otimes x_2)\\
(u_1\ast u_2)(x_1\ast x_2)&=-\chi\cdot u_1(x_1)\otimes u_2(x_2)=
\end{aligned}\right.
\]

\margin{$\uparrow$ signe étrange !}

\[
\begin{aligned}
(x_1\unsure{+|\ast}u_1)\ast(x_2+u_2)
&=x_1\ast x_{\unsure{1|2}})+
\underbrace{x_1\ast u_2}_{\underbrace{\pi_1(x_1)}_{b_1}\otimes u_2}
+\underbrace{u_1\ast x_2}_{+\,u_1\otimes\underbrace{\pi_2(x_2)}_{b_2}}
+\underbrace{u_1\ast u_2}_{\chi u_1\otimes u_2}
\end{aligned}
\]

% archive sheet number: 135
\scanpage{137}

\mspage{471}

qui a imprégné ma première enfance et qui était "mienne" déjà à ma naissance. Divers signes m'ont fait soupçonner plus d'une fois que cette note est "yin", que ce sont les qualités "féminines" qui dominent dans mon être, quand celui-ci trouve occasion de se manifester spontanément, en les instants où il est libre des conditionnements de toutes sortes qui se sont accumulés en moi depuis l'enfance. Pour le dire autrement: il se pourrait que ce qui est force créatrice en mon corps et en mon esprit, ce que j'ai appelé parfois "l'enfant" ou "l'ouvrier" en moi (par opposition au "patron" qui représente la structure du moi, c'est à dire ce qui est conditionné en moi, la somme ou le résultat du conditionnement accumulé en ma personne) - que cette force soit plus "féminine" encore que "virile" (alors que par nature et nécessité elle est l'un, et l'autre).

Ce n'est pas le lieu ici de passer en revue tous ces "signes". La chose importante d'ailleurs n'est \mbox{p a s} si cette note dominante profonde en moi est "féminine", ou si elle est "virile". C'est plutôt, que je sache en chaque instant \mbox{ê t r e}\quad\mbox{m o i - m ê m e}, en accueillant sans réticence aussi bien les traits et les pulsions en moi par lesquels je\unsure{:} suis "femme", que ceux par lesquels je suis "homme", et en leur permettant de s'exprimer librement.

Quand j'étais enfant, dans ces premières années, il n'était pas rare que \unsure{de|des} personnes étrangères me prennent pour une fille - sans d'ailleurs que la chose crée jamais en moi le moindre malaise, le moindre sentiment d'insécurité. \unsure{C'étai|C'était} surtout ma voix je crois qui faisait cet effet, une voix très claire, aigüe - sans compter que j'avais de longs cheveux (le plus souvent ébouriffés), peut-être simplement parce que ma mère (qui ne manquait pas d'autres soucis) ne prenait pas souvent le loisir de me les couper tant soit peu. J'étais par ailleurs fort comme un turc et les jeux un peu violents ou casse-cou ne me déplaisaient pas, ce qui n'empêchait nullement un penchant pour le silence, voire même pour la solitude, et un penchant également pour jouer à la poupée (*). Je n'ai pas souvenir que \unsure{quelqu'un|quèlqu'un} se soit moqué de moi à ce propos, mais la chose sûrement n'a guère pu manquer de se produire ici et là. Si de tels incidents ont passé sans laisser trace de blessure ou d'humiliation, c'est sûrement qu'ils ne recueillaient aucun écho ou amplification, par quelque sentiment d'insécurité en moi, alors que l'acceptation de celui que j'étais, par ceux qui seuls pour moi comptaient vraiment, était au delà de toute question. La moquerie n'aurait pu m'atteindre, elle ne pouvait que se retourner contre celui qui devait m'apparaître comme bien sot, pour faire mine de trouver à redire à la chose la plus naturelle du monde.

\noindent\rule{5em}{0.4pt}\par
(*) Si ce penchant paraît rare chez les petits garçons, c'est surtout je crois parce qu'il est systématiquement découragé par l'entourage.

% archive sheet number: 136
\scanpage{138}

\[
\deleted{(f_1\ast f_2):(\unsure{x}_1\ast x_2)=}
\qquad (f_1\ast f_2)\mathord{\upharpoonright} E_1\otimes E_2:
 E_1\otimes E_2 - L_1\otimes L_2
\]

\[
\boxed{\begin{aligned}
 f_1\ast f_2(x_1\ast x_2)&=\lambda_1\pi_1(x_1)\otimes f_2(x_2)
 + f_1(x_1)\otimes\lambda_2\pi_2(x_2)^{\deleted{\unsure{\lambda_2+}}}
 -\unsure{\chi|\lambda}f_1(x_1)\otimes f_{\unsure{1|2}}(x_{\unsure{1|2}})\\
 \text{i.e.}\quad f_1\ast f_2\mathord{\upharpoonright} E_1\otimes E_{\unsure{1|2}}\quad
 [E_1\otimes E_2\unsure{-|\to} L_1\otimes L_2]&\overset{\unsure{\prime|\cdot}}{=}\lambda_1\pi_1\otimes f_2+
 f_1\otimes\lambda_2\pi_2-\chi f_1\otimes f_2
\end{aligned}}
\]

\[
 f_1\ast f_2(u_1\otimes u_2)=\lambda_1\lambda_2(u_1\otimes u_2)
\]

\deleted{\unsure{Si|$f_1$}\quad $x$}\qquad Vérifions

\[
\begin{aligned}
 f_1\ast f_2(x_1\ast u_2)&=\lambda_1\lambda_2\pi_1(x_1)\otimes u_2\\
 \lambda_1\pi_1(x_1)\unsure{)}\otimes\lambda_2 u_2
 &+\deleted{f_1(x_1)\otimes\underline{\lambda_2\chi}\,u_2}
 -\deleted{\underline{\chi f_1(x_1)\otimes\lambda_2}\,u_2}\qquad\text{OK}
\end{aligned}
\]

\[
\begin{aligned}
(\pi_1\ast\pi_2)(x_1\ast x_2)&=\chi\,\big\{\pi_1(x_1)\otimes\pi_2(x_2),\\
(\pi_1\ast\pi_2)(u_1\otimes u_2)&=\chi^2\deleted{\unsure{\pi}}u_1\otimes u_2
\end{aligned}
\qquad\left|\quad \frac1\chi\pi_1\unsure{\ast|\otimes}\pi_2=\pi\right.
\]

\[
 x_1\ast x_2\unsure{\equiv|=}\unsure{\psi_1(x_1)|\varphi_1(x_1)|\deleted{\psi_1(x_1)}}
\]

\diagram{$\begin{array}{cc}L_i&p_i\\ M_i&q_i\end{array}$}

\[
 1+\underline{p_1q_1},\qquad 1+\underline{p_2q_2}\qquad 1+(p_1p_2)(q_1q_2)
\]

\[
 \deleted{[x,y]=}\qquad x\to\deleted{\unsure{p|\psi}(x)-\unsure{xf(}}\psi(x)T-\chi x
\]

% archive sheet number: 137
\scanpage{139}

\mspage{472}

Je savais bien d'ailleurs que ce genre de sottise un peu étrange n'est nullement chose rare, que la seule vue de la nudité peut être cause de scandale! Pourtant aussi loin que je pouvais me souvenir, j'avais eu l'occasion de voir ma mère, mon père et ma soeur nus, et toute occasion aussi de satisfaire ma légitime curiosité quant à comment chacun d'eux ainsi que moi-même étions faits. Il était bien évident qu'il n'y avait nulle cause de scandale dans la conformation des hommes ou des femmes, qui me paraissait décidément très bien comme elle était - et plus particulièrement (je n'en faisais nullement mystère) celle des femmes.

(\textsuperscript{108}) (5 octobre) C'est en 1933, alors que j'étais dans ma sixième année, que se place le premier tournant crucial dans ma vie, qui a été en même temps un tournant crucial aussi dans la vie de ma mère comme de mon père, dans leur relation l'un à l'autre comme dans celle à leur enfants. C'est l'épisode de la \unsure{destruction|destructio} violente et définitive de la famille que nous formions tous les quatre, destruction dont j'ai été le premier et le seul, quarante-six ans \unsure{.}plus tard, à faire le constat et à suivre les péripéties, dans la correspondance de mes parents et dans un ou deux souvenirs exsangues, énigmatiques et tenaces, patiemment sondés et déchiffrés - bien longtemps après la mort de mon père et celle de ma mère (*).

Ce n'est pas mon propos de m'étendre ici sur ce que j'ai appris et compris au cours de ce long travail, au sujet de la portée et du sens de cet épisode. J'ai déjà fait allusion il y a trois jours à ce tournant (**), comme marquant la fin brutale de la première des trois grandes périodes, dans l'histoire des épousailles du yin et du yang en moi. En décembre 1933, je me trouve largué en toute hâte dans une famille étrangère, que moi, ni ma mère qui m'y amenait depuis Berlin, n'avions jamais vue. En fait, ces gens inconnus chez qui elle m'amenait étaient simplement les premiers venus qui veuillent bien de moi comme "pensionnaire" pour une pension plus que modique, et avec aucune garantie d'aucune sorte que celle-ci serait jamais payée, alors que ma mère s'apprêtait à rejoindre au plus vite mon père, qui se morfondait à l'attendre à Paris. C'était une chose en-

\noindent\rule{5em}{0.4pt}\par
(*) Mon père est mort à Auschwitz en 1942, ma mère est morte en 1957. Le travail dont je parle ici s'est poursuivi entre août 1979 et octobre 1980.

(**) Voir fin de la note "Yang enterre yin - ou le muscle et la tripe", n° 106.

% archive sheet number: 138
\scanpage{140}

\diagram{\begin{tikzcd}[ampersand replacement=\&,row sep=small]
 M_1\arrow[d,bend left,"{q_1}"]\\
 E_1\arrow[d,"{\alpha_1}"']\arrow[d,bend left,"{p_1}"]\\L_1
\end{tikzcd}}

\diagram{\begin{tikzcd}[ampersand replacement=\&,column sep=small]
 \&\&\deleted{M_1\otimes M_2}\&\\
 (E_1\otimes L_2)\oplus(L_1\otimes E_2)\arrow[r]\arrow[dr,"{\pi}"']\&
 {E_1\otimes E_2\atop x_1\otimes x_2}\arrow[r,dash]\arrow[dr,"{\unsure{\lambda_1}}"]\&
 {M_1\otimes M_2\atop m_1\otimes m_2}\arrow[r,dash]\&0\\
 \& L_1\otimes L_2\arrow[r,dash]\&
 {E\atop x_1\ast x_2}\arrow[u,phantom,"{\deleted{\unsure{\nearrow|\wr}}}"]\arrow[r]\arrow[l,bend left,"p"]\&
 {\deleted{\unsure{L}}M_1\otimes M_2\atop m_1\otimes m_2}\arrow[l,bend left,"q"]
\end{tikzcd}}

\[
\text{\unsure{cocart|cocart.}}\qquad x_1\ast x_2^{\unsure{7|]}}\longmapsto m_1\otimes m_2
\]

\[
\begin{aligned}
q(m_1\otimes m_2)&=q_1(m_1)\ast q_{\unsure{1|2}}(m_2)\\
p(\deleted{\unsure{m}}x_1\deleted{\unsure{\otimes}}\ast x_2)&=x_1\ast x_2-q(m_1\otimes m_2)\\
 &=x_1\ast x_2-q_{\unsure{1}}(m_1)\ast q_2(m_2)\\
q(m_1)&=x_1-p_1(x_1)\\
q(m_2)&=x_2-p_2(x_2)
\end{aligned}
\]

\[
\begin{aligned}
p(x_1\ast x_2)&=x_1\ast x_2-(x_1-p_1(x_1))\ast(x_2-p_2(x_2))\\
 &=\unsure{+}x_1\ast p_2(x_2)+p_{1}(x_1)\ast x_2-p_1\underline{(x_1)\ast p_{\unsure{2}}(x_2)}\\
 &\deleted{=\unsure{\chi|\lambda}(x_1\otimes p_2(x_2)+p_1(x_1)\otimes x_2-\unsure{\chi}p_1(x_1)\otimes p_2(x_{\unsure{1|2}}))}
\end{aligned}
\]

\[
\begin{aligned}
u_1\ast u_2&=\chi u_1\otimes u_2\\
\deleted{p(x_1\ast{}}\qquad p(u_1\ast u_2)&=\chi u_1\otimes u_2
\end{aligned}
\]

\[
\boxed{(f_1\ast f_2)(x_1\ast x_2)=
\underbrace{\lambda_1x_1}_{\lambda_1\pi_1(x_1)}\ast f_2(x_2)
+f_1(x_1)\ast\lambda_2 x_2-\chi f_1(x_1)\otimes f_2(x_2)}
\]

\deleted{Si $x_1,x_2\in\unsure{L}$}

\[
\begin{aligned}
\text{Si }\quad x_1&\deleted{\in}=u_1\in L_1\\
x_2&=u_2\in L_2
\end{aligned}
\]

\[
\begin{aligned}
2\lambda_1\deleted{\unsure{\lambda_2}}\chi x_1\mathop{\otimes}_{\lambda_2}x_2
 &-\deleted{\unsure{2}}\chi\lambda_1\lambda_2 x_1\otimes x_2\\
 &=\\
\deleted{\lambda\chi}\qquad&\deleted{(2-\chi)\lambda_1\lambda_2 u_1\otimes u_2}\\
&2-\chi=1
\end{aligned}
\]

\[
\underbrace{xy-(x-x)(y-y)}_{xy+yx-xy}=xy\qquad\Big/
\]

\[
\begin{aligned}
\deleted{\unsure{f}}\qquad (f_1\ast f_2)(\unsure{x_1|u_1}\otimes u_2)&=\lambda_1\lambda_2 u_1\otimes u_2\\
 &=f_1(u_1)\otimes f_2(u_2)
\end{aligned}
\]

% archive sheet number: 139
\scanpage{141}

\mspage{473}

tendue entre mes parents que tout allait être pour le mieux tant pour moi à Blankenese (près de Hambourg), que pour ma soeur qui depuis quelques mois avait été larguée à la fin des fins dans une institution à Berlin pour enfants handicapés (où on avait bien voulu d'elle, bien qu'elle n'était pas plus handicapée que moi ou nos parents).

En dénouement à six mois étranges, lourds de menace sourde et d'angoisse, je me suis retrouvé du jour au lendemain dans un monde totalement différent du seul monde que j'avais connu dans ma vie, celui formé par mes parents et par ma soeur et moi. Je m'y retrouvais comme un parmi un groupe de pensionnaires, qui mangions à part de la famille et faisions figure d'enfants de deuxième catégorie pour les enfants de la maison, lesquels formaient un monde à part et nous regardaient de haut. De ma mère je recevais une lettre hâtive et guindée de loin en loin, et de mon père jamais une ligne de sa main, pendant les cinq ans que j'y suis resté (jusqu'en 1939, à la veille de la guerre, quand j'ai fini par \unsure{rejoind|rejoindre} mes parents sous la pression des évènements).

Le couple qui m'avait accueilli m'a vite pris en affection. Aussi bien lui ancien pasteur qui avait quitté le sacerdoce et vivait d'une maigre pension et de leçons particulières de latin, de grec et de mathématiques, que sa femme pétillante de vie et parfois de malice, étaient des gens peu ordinaires, \unsure{attachant|attachants} à bien des égards. Lui était un humaniste de vaste culture qui s'était un peu égaré dans la politique, et avait eu maille à partir avec le régime nazi, qui a fini par le laisser tranquille. Après la guerre j'ai renoué et je suis resté en relations suivies avec eux jusqu'à la mort de l'un et de l'autre (*).

De lui et surtout d'elle, tout comme de mes parents, j'ai reçu du meilleur comme aussi du pire. Aujourd'hui, avec un long recul, je leur suis reconnaissant (comme je le suis à mes parents) pour ce "meilleur", comme aussi pour ce "pire". C'est ce meilleur et ce pire que j'ai reçu, de mes parents d'abord, d'eux \unsure{ensuite|ensuit} qui a formé le plus gros du volumineux "paquet" que j'ai reçu en partage dans \unsure{mo|mon} enfance (comme chacun reçoit le sien...)\handwritten{\unsure{,}}\unsure{,}qu'il m'appartenait de déballer et d'\unsure{exminer|examiner}. Ils font partie de la substance, de la richesse de mon passé, dont il ne tient qu'à moi de nourrir mon présent.

Mon nouveau milieu était tout ce qu'il y a de "comme il faut" et de conformiste à beaucoup d'égards, avec en tous cas les attitudes répressives de rigueur pour tout ce qui concerne le corps et, plus particulièrement, le sexe. Il a \unsure{fall|fallu}

\noindent\rule{5em}{0.4pt}\par
(*) Elle est morte à l'âge de 99 ans, il y a deux ans, et j'ai encore pu la voir morte, en tête à tête avec elle, la veille de l'enterrement.

% archive sheet number: 140
\scanpage{142}

Si on a pu définir $F_{11}$, \unsure{et|avec} commutativité des deux autres carrés, alors l'hom $F$ est un hom d'ext. au sens strict (induit $\mathrm{id}$ sur $L_1\otimes L_2$ et sur $M_{\unsure{2|1}}\otimes M_2$), donc \unsure{est} un iso.

Contraintes sur $F_{11}:E_1\otimes E_2 - E_1\otimes E_2$ \unsure{endom.}

1°) Connu sur $E_1\otimes_0 E_2$, c'est le \unsure{composé|composée}:
$E_1\otimes_0 E_2\xrightarrow{\pi}L_1\otimes L_2\longleftarrow E_1\otimes E_2$
($\chi$-rétraction de $E_1\otimes_0 E_2$ sur le sous-module $L_1\otimes L_2$)

i.e. $F_{11}\circ\mu$ est connu

\diagram{\begin{tikzcd}[ampersand replacement=\&,column sep=small]
E_1\otimes_0\underline E_2\arrow[r,"{\mu}"]\arrow[dr,dashed,"{\pi}"']\&
\deleted{\unsure{L}}E_1\otimes E_2\arrow[r,"{F_{11}}"]\&
E_1\otimes E_2\\
\& L_1\otimes L_2\arrow[ur,dashed,"{\alpha_1\otimes\alpha_2}"']\&
\end{tikzcd}}

2°) $\unsure{\nu|v}\circ F_{11}$ est connu, de façon \unsure{précise} on a

\diagram{\begin{tikzcd}[ampersand replacement=\&,column sep=small]
E_1\otimes E_2\arrow[r,"{F_{11}}"]\arrow[dr,"{\psi_1\otimes\psi_2}"']\&
E_1\otimes E_2\arrow[r,"{\nu}"]\&
\deleted{\unsure{E}}E_1\otimes^{0}E_2\\
\&M_1\otimes M_2\arrow[ur,"{\varpi}"']\&
\end{tikzcd}}

\margin{\deleted{\uline{NB} $\pi\mathord{\upharpoonright}(\unsure{\alpha_1}\otimes L_2)=\chi\,\mathrm{id}_{L_1\otimes L_2}$}
NB\quad $\pi\circ(\alpha_1\otimes\alpha_2)=\chi\,\mathrm{id}_{L_1\otimes\unsure{L_2}}$
$(\psi_1\otimes\psi_2)\circ\varpi=\chi\,\mathrm{id}_{M_1\otimes\unsure{M_2}}$}

\[
\boxed{\begin{aligned}
 &\chi\,\mathrm{id}-(\pi'_1+\pi'_2)+\gamma\,\pi'_1\otimes\pi'_2\\
 &\gamma\chi=2\qquad \text{e.g.}\quad\chi=2,\ \gamma=1
\end{aligned}}
\]

dans le cas $\chi=2$, on aura

\[
\left\{\begin{aligned}
 F_{11}(\deleted{\unsure{x_1\otimes{}}}x)&=\bar x+x\\
 \bar x_i&=x_i-\pi_i(x_i)
\end{aligned}\right.
\qquad
\begin{aligned}
\bar{x}_1\otimes\bar{x}_2&=\bar{x}_1\unsure{\ominus|\otimes}\bar{x}_2\\
\text{i.e.}\quad &\boxed{\pi_i(x_i)=x_i-\bar x_i}
\end{aligned}
\]

\margin{NB La \unsure{donnée} de $\pi_1$ \unsure{équivaut} à celle de $\sigma_1$ \unsure{involution} qui \unsure{induisant|induit} $-\mathrm{id}_{L_1}$ et $+\mathrm{id}_{M_1}$}

% archive sheet number: 141
\scanpage{143}

\mspage{480}

Dans l'optique de la réflexion de ces derniers jours, c'est ce moment de retrouvailles avec mon enfance, crue perdue et morte une longue vie durant, qui marque la fin de la "deuxième période" de mon itinéraire spirituel: celle de la prédaminance, dans ma vie personelle, des \mbox{m é c a n i s m e s}\quad\mbox{é g o t i q u e s}, à l'encontre des forces créatrices, des forces de connaissance et de renouvellement, qui avaient passé par une stagnation presque complète de quarante ans.C'est aussi l'époque de la prépondérance d'une "certaine force", d'une force à caractère presque exclusivement "viril", à l'image des valeurs en \unsure{honneu|honneur} dans le monde environnant, aux dépens des aspects et forces profondes "féminins" de mon être, ignorés et refoulés (avec un succés jamais complet, Dieu merci!).

La toute première intuition sur la nature destructrice de cette force-là, laquelle avait dominé ma vie comme celle de ma mère, et celle d'autres femmes aussi qui avaient été importantes dans ma vie - cette intuition fait une courte apparition en ces jours de maturation intense, à la faveur sûrement du ressurgissement de l'énergie yin, "féminine", dans mon appréhension consciente des choses Contrairement à ce dont je croyais hâtivement me souvenir tantôt, cette apparition n'a pas lieu dans la méditation de la veille des retrouvailles, mais quelques heures après celles-ci, dans une courte méditation sur le sens de ce qui venait de se passer. L'intuition naît et prend forme tout à la\unsure{.}fin des quelques pages de notes de cette méditation. Je perçois la nature destructrice de \unsure{cete|cette} "force" (qu'aujourd'hui j'appellerais "force superyang", c'est à dire à dominante yang excessive) chez ma mère d'abord, puis chez d'autres femmes, pour \unsure{enchainer|erchainer} avec ces lignes finales:

\begin{quote}
" Quant à la "force" en moi-même, c'est elle certes qui a fait de moi la cible et l'objet, attendus une jeune vie durant, de la haîne secrète et du ressentiment de M., puis de J., puis de S. - d'une haîne déposée en elles longtemps avant qu'elles ne connaissent mon existence, en les jours désemparés d'une enfance privée d'amour."
\end{quote}

Le mot "enfance", dans la dernière ligne encore qui témoigne d'un jour important entre tous dans ma vie, y apparaît d'ailleurs pour la dernière fois \unsure{pour|pou} près de trois ans! Quant à l'intuition sur la nature de la force superyang en moi, comme provocatrice de réactions antagonistes, voire de haîne et de ressentiment, elle a eu tendance (il me semble) à sombrer un peu dans l'oubli jusqu'à ces tout derniers jours encore. Plus précisément, elle est restée présente seulement dans ma perception de certaines relations importantes dans ma vie (et \unsure{sutout|surtout}, des relations avec des femmes que j'ai aimées). Par contre, elle n'a \unsure{guère|guèr}

% archive sheet number: 142
\scanpage{144}

\[
E_1\otimes E_2^{\deleted{\unsure{\ast}}}\qquad
E_1\mathbin{\unsure{\vee|\ast|\circledast}}E_2=
E_1\otimes E_2\underset{(E_1\otimes_0 E_2)}{\oplus} L_1\otimes L_2
\]

$\deleted{\unsure{f|u}}u_i$ endom de $E_i$ tq. $u_i(L_i)\subset L_i$
\quad $u_i\mathord{\upharpoonright}L_i=u^{\unsure{\prime|!}}_i:L_i - L_i$

\deleted{$u_2$ endom de $E_2$}

\[
\begin{aligned}
\deleted{u_1}f_1\otimes\deleted{u_2}f_2&\quad\text{sur }E_1\otimes E_2\\
\deleted{u'_1}f^{!}_1\otimes\deleted{u'_2}f^{!}_2&\quad -\quad L_1\otimes L_2
\end{aligned}
\]

\[
\begin{aligned}
\deleted{u_1\otimes u_2}\qquad \deleted{u}f(x_1\ast x_2)&=\deleted{u_1}f_1(x_1)\ast\deleted{u_2}f_2(x_2)\\
\deleted{u}f(u_1\otimes u_2)&=f^{!}_1(u_1)\otimes f^{!}_{\unsure{2}}(u_2)
\end{aligned}
\]

\diagram{\begin{tikzcd}[ampersand replacement=\&,column sep=small,row sep=small]
 E_1\otimes L_2\arrow[d]\& {E_1\otimes E_{\unsure{1|2}}\atop x_1\otimes u_2}
 \arrow[r]\& f_1(x_1)\ast f_2(u_2)\\
 x_1\otimes u_2\arrow[ur,maps to]\arrow[dr,maps to]\&\&\\
 \&\pi_1(x_1)\otimes u_2\arrow[r]\& f^{!}_1(\pi_1(x_1))\otimes f^{!}_2(u_2)
\end{tikzcd}}

\[
 f_1(x_1)\ast\unsure{v_2|w_2|u_2}=f^{!}_1(\pi_1(x_1))\otimes v_2
\]

\[
\begin{aligned}
 \deleted{f(x}\quad &\text{\uline{Ex}}\quad f_1=\\
 f_1(x_1)&=f_1(\unsure{T|\pi}\qquad \deleted{\unsure{f_2}}\\
 x_1&-\pi_1(x_1)
\end{aligned}
\]

\diagram{\begin{tikzcd}[ampersand replacement=\&]
E_1\otimes E_2\arrow[r]\& \deleted{\unsure{O|E}}L_1\otimes L_2\\[-2pt]
x\& 
\end{tikzcd}}

\diagram{\begin{tikzcd}[ampersand replacement=\&,row sep=small]
 E_1\arrow[d,bend left,"f_1"']\&E_2\arrow[d,bend left,"f_2"]\\L_1\&L_2\\
 \lambda_1\unsure{)}\&\lambda_2
\end{tikzcd}}

\diagram{\begin{tikzcd}[ampersand replacement=\&,column sep=small]
 \&E_1\otimes E_2\arrow[r,"{\chi\,\deleted{\pi_1}f_1\otimes\deleted{\pi_2}f_2}"]\&L_1\otimes L_2\\
 E_1\otimes L_2\arrow[ur]\arrow[dr]\&\&\\
 \&L_1\otimes L_2\arrow[uur,"{\deleted{\unsure{\chi_1\chi_2}}\lambda_1\lambda_2}"']\&
\end{tikzcd}}

\[
\begin{aligned}
 f_1&\qquad f_2\\ \lambda_1&\qquad\lambda_2\\
 \frac{\lambda_{\unsure{2|1}}}{\lambda_{\unsure{12|2}}}\,f_2&\qquad
 \chi f_1(x_1)\otimes\overbrace{f_2(u_2)}^{\lambda_2u_2}\\
 &\qquad\deleted{f_1(\unsure{\pi_1|f_1}(x_1))\otimes u_2}\\
 &\qquad \lambda_1\deleted{\unsure{\chi|\lambda_2}}\lambda_2\pi_1(x_1)\otimes\deleted{\unsure{\lambda_2}}u_2
\end{aligned}
\]

\[
\begin{aligned}
 \deleted{f_1(}\quad x_1\otimes u_2&\qquad\chi(\pi_1(x_1)\otimes u_2)\\
 \deleted{f_1\otimes f_2}&\qquad \deleted{\lambda_1\lambda_2(\pi_1(x_1)\otimes u_2)}\\
 &\qquad f_1(\pi_1(\unsure{x'_1|x_1}))\otimes f_2(u_2)\\
 &\qquad\deleted{\lambda_2 f_1(x_1)=\lambda}\\
 &\qquad \lambda_1\pi_1(x_1)=f_1(x_1)\chi
\end{aligned}
\]

% archive sheet number: 143
\scanpage{145}

\mspage{474}

pourtant plusieurs années, je crois, avant que je m'intériorise et ne reprenne à mon compte ces attitudes là, comme la honte de me montrer nu, allant de pair \unsure{ave|avec} une relation ambigüe avec mon corps. Cette honte,inculquée dès le jeune âge, est un des aspects d'une division profonde, où le corps est objet d'un tacite \unsure{mépris|mépri} alors que les valeurs dites "culturelles" (confondues avec des capacité intellectuelles de mémorisation et autres) sont montées en épingle. Cette division en m\handwritten{oi} est restée igorée jusqu'à ma quarante-huitième année, où elle a commencé à se r\handwritten{é}soudre. C'est là le deuxième grand tournant de ma vie, qui marque l'avènement de la "troisième période" dans l'histoire de ma relation à moi-même, c'est-à-dire \unsure{asi|aussi} celle de ma relation à mon corps, et à "l'homme" et à "la femme" en moi. Mais auparavant j'ai eu ample occasion de contribuer à transmettre cette division à mes enfants (*)\handwritten{,} que j'ai pu voir la transmettre à leur tour...

J'ai fait allusion déjà hier (**) au "basculement" qui a fini par avoir lieu en moi. Avec un décalage de plus de deux ans après l'arrachement au milieu familial initial (ou pour mieux dire, après la \quad\mbox{d e s t r u c t i o n}\quad de ce milieu), ce basculement consacre la mise en place des mécanismes répressifs courants, dont mon enfance avait eu la rare chance d'être exempte jusque là. J'ai détecté jusqu'à présent deux grandes forces de nature répressive, qui ont dominé ma vie d'adulte et une grande partie de mon enfance ($^{108}1$). Je crois pouvoir dir\handwritten{e} que leur apparition ne s'est pas faite progressivement, mais que dans mon cas ce\handwritten{s} mécanismes sont apparus plus ou moins du jour au lendemain et dans toute leur force, comme conséquence d'un \quad\mbox{c h o i x}\quad délibéré, au niveau inconscient. J'a\handwritten{i} qualifié précédemment ce choix d'"abdication", mais en même temps c'était aussi un puissant principe d'action: le "je serai comme "eux"" (et pas "comme moi") signifiait aussi: je vais "miser" sur "la tête", pas plus mauvaise chez moi que chez quiconque après tout, et me battre et "les" battre avec leurs propres armes

L'un de ces mécanismes, et celui qui m'interesse surtout ici, est un des plus communs qui soit: c'est la \quad\mbox{r é p r e s s i o n}\quad\mbox{d e}\quad\mbox{m e s}\quad\mbox{\unsure{t r a i t|t r a i t s}}\quad "\mbox{f é m i n i n s}" (ou ceux ressentis comme tels par les consensus courants), au profit de valeurs "viriles". L'endroit de la médaille était bien sûr l'inverstissement à fond sur mes traits et aptitudes ressentis comme "virils" et le développement à outrance de ceux-ci, qui ont pris une place démesurée.

\noindent\rule{.21\linewidth}{.35pt}

(*) Du moins, aux quatre parmi eux que j'ai contribué à élever. Le cinquième et dernier est élevé par sa mère, et jusqu'à présent il ne s'est pas présenté une occasion propice pour seulement faire connaissance, lui et moi.

(**) Voir les débuts de la note précédente "Eclosion de la force - ou les épousailles", note n° 10\deleted{\unsure{8|9}}\handwritten{\unsure{7|9}}.

% archive sheet number: 144
